% Folder 138, batch 2 (archivists' pages 21-40).
% Pass: Opus 5.5 (claude-opus-5-5), 2026-09-30 — first pass, unchecked against the pages by a human.
%
% Page 23 carries only show-through from its neighbour and is skipped.
% Notation fixed for the batch: \mathfrak{S}_3 for his S_3; T for
% \varprojlim \mu_n; \Pi_n, \Pi_\infty; \Gamma = Gal(\bar Q/Q); his
% "(G).½ H" (semi-direct product) set as G \cdot_{1/2} H.
\documentclass[11pt,a4paper]{article}
\input{../preamble/grothendieck}

\folder{138}
\batch{2}
\pages{21}{40}
\foldertitle{Cartes. Etude arithmétique (1976, 77 ?) : bon de commande (s.d.), notes manuscrites (s.d.).}
\dating{[vers 1976-1978]}
\watermark{Édition de démonstration}

\begin{document}

\page{21}
\note{feuillet de calculs épars, sans phrase suivie ; transcrits dans l'ordre de la page, de haut en bas}

$\Gamma$ opère sur $\overline{X}'_{\xi}$ \quad ($\xi \in \mathbb{Q}^{*}$)

$\Gamma$ opère trivialement sur $\overline{X}'_{-1}$ \quad
$\bigl[\,\overline{X}'_{(-1)^{1/n}}$, \ $\overline{X}'_{\xi^{1/n}}\,\bigr]$
\note{les indices des deux derniers $\overline{X}'$ sont lus sans assurance}

\[
  \exp\frac{2i\pi}{2n} = \exp\frac{i\pi}{n}
\]

\[
  \mathbb{Q}^{*} \longrightarrow Z^{1}(\Gamma, \widehat{\mathbb{Z}}^{*}) \longrightarrow H^{1}(\Gamma, \widehat{\mathbb{Z}})
\]
\note{devant $Z^{1}$, « $\mathrm{Hom}($ » biffé}
\[
  \mathbb{Q}^{*} \simeq \mathbb{Z}/2\mathbb{Z} \oplus \mathbb{Z}^{(\mathbb{P})}, \qquad
  \varphi(\gamma\gamma') = \varphi(\gamma) + \gamma\varphi(\gamma'), \qquad
  \varphi(\gamma) = a - \gamma a
\]
\[
  \mathbb{Q}^{*} \longrightarrow Z^{1}(\Gamma, \widehat{\mathbb{Z}}_{\ell}), \qquad
  \widehat{\mathbb{Z}}^{\mathbb{P}} \longrightarrow Z^{1}, \qquad
  \varprojlim_{n} Z^{1}(\Gamma, \mathbb{Z}/n\mathbb{Z})
\]
\note{un « $\mathbb{Z}[$ » inachevé sous $\mathbb{Z}/2\mathbb{Z}\oplus\mathbb{Z}^{(\mathbb{P})}$, et un griffonnage biffé ($\mathbb{Q}^{*}$, $\mathbb{Z}^{\mathbb{P}}$, $\mathbb{Z}_{\ell}$) au-dessus de $\widehat{\mathbb{Z}}^{\mathbb{P}}$}

\[
  \mathbb{Q}^{*}/\mathbb{Q}^{*n} \simeq H^{1}(\Gamma, \mu_n) \simeq Z^{1}(\Gamma, \mathbb{Z}/n\mathbb{Z})
\]
\struck{$a \in \mathbb{Q}^{*}$ \quad $\overline{\mathbb{Q}}$}

$(a, \ \xi) \in \overline{\mathbb{Q}}$, \quad $\xi^{n} \in \mathbb{Q}^{*}$, \quad
$\mathbb{Q}^{*} \subset \overline{\mathbb{Q}}^{*}$, \quad ${}_{n}(\overline{\mathbb{Q}}^{*}/\mathbb{Q}^{*})$

$\sqrt[n]{p_1}, \ldots, \sqrt[n]{p_\nu}$ \qquad
$K = K_{A,n} = \mathbb{Q}(\sqrt[n]{p_1}, \ldots, \sqrt[n]{p_\nu})$

$\mathrm{T}$ \quad $K^{*}/\mathbb{Q}^{*}$

\page{22}
\note{feuillet de calculs épars}

\[
  \mathbb{G}_{m,\mathbb{Q}} \xrightarrow{\ \lambda \mapsto \lambda^{n}\ } \mathbb{G}_{m,\mathbb{Q}}
  \qquad\qquad \mathbb{Q} \longrightarrow \overline{\mathbb{Q}}
\]
\note{écrit verticalement sur la page : la flèche $\lambda\mapsto\lambda^n$ descend de $\mathbb{G}_{m,\mathbb{Q}}$ à $\mathbb{G}_{m,\mathbb{Q}}$, puis un trait vers $\mathbb{Q}\to\overline{\mathbb{Q}}$}

\[
  1 \longrightarrow \mu_n \longrightarrow \mathcal{E}'_n \rightleftarrows \Gamma' \longrightarrow 1
\]
\[
  \mathcal{E}'_n = \bigl[(\mathbb{Z}/6n\mathbb{Z})^{*} \cdot \mu_\nu\bigr] \cdot \mu_n
  \qquad \bigl[(\mathbb{Z}/6n\mathbb{Z})^{*} \cdot \mu_\nu \to \mathrm{Aff}(\mu_n)\bigr]
\]
\struck{$\Gamma' \simeq (\mathbb{Z}/6n\mathbb{Z})^{*} \cdot \mu_\nu$}

\begin{tikzcd}[column sep=small, row sep=small, nodes={font=\scriptsize}]
1 \arrow[r] & \mu_\nu \arrow[r] & \Gamma'_\nu \arrow[r] & (\mathbb{Z}/\nu\mathbb{Z})^{*} \arrow[r] & 1 \\
1 \arrow[r] & \mu_\nu \arrow[r] \arrow[u, no head, "\Vert" description] & \Gamma' \arrow[r] \arrow[u] & (\mathbb{Z}/6n\mathbb{Z})^{*} \arrow[r] \arrow[u] & 1
\end{tikzcd}

\marginal{splitté, splitting dépend du choix de $\sqrt[\nu]{2}$ sur $K_\nu$.}
\note{la flèche de cette note pointe vers $\Gamma'_\nu$ ; l'indice de $\Gamma'$ en bas est illisible, lu $\Gamma'$ sans indice ; l'égalité verticale de gauche est tracée par un double trait}

\[
  \begin{array}{cccccl}
    E^{+} & R' & R & S_1 & S_0 & \\
    \wr & \wr & \wr & \wr & \wr & \bigr\}\ z \mapsto \tfrac{1}{z} \\
    E^{-} & R'' & R & S_1 & S_2 &
  \end{array}
\]

\[
  1 \longrightarrow \mu_\infty(\overline{\mathbb{Q}}) \longrightarrow \mathcal{E}_\infty \longrightarrow \Gamma \longrightarrow 1
\]
$\Gamma \cdot \mu_\infty(\overline{\mathbb{Q}})$ \quad [$\mu_\infty$ ($\simeq \widehat{\mathbb{Z}}$, non can.)]
\quad $\Gamma$ opère sur $\widehat{\mathbb{Z}}^{*}$

\struck{$\mathbb{Z}$, $R'$, $R''$, $1$}

$E^{+}, E^{-}, R, R', R'', S_1$ sont des torseurs sous $\mu_\infty$

\struck{$E^{+} \simeq (E^{-})^{\vee}$, $R' \simeq$}

$E^{+} \wedge E^{-} = 1$, \ $R' \wedge R'' = 1$, \ $R^{2} = 1$

$E^{+\,6} = 1$

% Page 23 is blank but for show-through.

\page{24}
\textbf{1.1} Soit $n \in \mathbb{N}^{*}$, et posons $C_n \subset \mathbb{P}^{2}_{\mathbb{Q}}$ équation projective
\[
  \text{(1)} \qquad x^{n} + y^{n} + z^{n} = 0 .
\]
Le groupe $\mu^{3}_{n\,\mathbb{Q}}$ opère sur $\mathbb{E}^{3}_{\mathbb{Q}}$ par matrices diagonales, donc posant
\[
  \text{(2)} \qquad \Pi_{n\,\mathbb{Q}} = \mu_n^{3}/\mu_n \ (\text{diag.}) \simeq \mu_n \otimes_{\mathbb{Z}} (\mathbb{Z}^{3}/\mathbb{Z})
\]
\marginal{NB on écrit $\mu_n$ au lieu de $\mu_{n\,\mathbb{Q}}$}
opère sur $\mathbb{P}^{2}_{\mathbb{Q}}$, en laissant $C_n$ invariant. La courbe $C_n/\Pi_{n\,\mathbb{Q}}$ s'identifie à la courbe \add{$C_1$} qui est donnée, en coordonnées homogènes, par
\[
  \text{(3)} \qquad X + Y + Z = 0
\]
qu'on identifie à $\mathbb{P}^{1}_{X}$, en envoyant les points de $C_1$ pour lesquels $X$, $Y$ ou $Z$ sont nuls resp.\ en $0, 1, \infty$. Donc écrivant $C_1$ en coordonnées affines, faisant $Z = 1$ i.e.\ $X/Z = \xi$, $Y/Z = \eta$, sous la forme
\[
  \text{(3')} \qquad \xi + \eta + 1 = 0
\]
i.e.
\[
  \eta = -1 - \xi ,
\]
$\xi = 0$ correspond à $\eta = -1$, $\eta = 0$ à $\xi = -1$, donc posant
\[
  \text{(3'')} \qquad t = -\xi \quad \text{d'où} \quad \xi = -t, \ \eta = t - 1, \ \zeta = 1
\]
les valeurs de $t$ pour lesquelles $X = 0$, $Y = 0$, $Z = 0$ sont resp.\ $t = 0, 1, \infty$.

$C_n$ est étale sur $C_1 = \mathbb{P}^{1}_{\mathbb{Q}}$ sauf en $0, 1, \infty$, i.e.\ c'est un revêtement principal $C_n^{*}$ de $C_1^{*} \simeq \mathbb{P}^{1}_{\mathbb{Q}} \setminus \{0, 1, \infty\}$, de groupe $\Pi_{n\,\mathbb{Q}}$.

\marginal{Il y a en $\alpha(i)$ ($i \in \{0, 1, \infty\}$) $n$ pts de ramification de $C_n$, dont \ill{} degré de ramification $n$. Comme cette \uncertain{variété}, elle est formée de $2n^{2}$ triangles, qui \ill{} \ill{} \ill{} \ill{} \ill{} \ill{} \ill{} fois\ldots}
\note{deux notes écrites le long de la marge gauche, en remontant ; le mot « variété » et la suite sont très incertains}

\page{25}
Passant à la limite, on trouve
\[
  \text{(5)} \qquad C_\infty = \varprojlim C_n, \qquad C_\infty^{*} = \varprojlim C_n^{*}
\]
sur lesquels opère
\[
  \text{(6)} \qquad \Pi_\infty = \varprojlim \Pi_n \simeq \mu^{\leftarrow}_\infty \otimes_{\widehat{\mathbb{Z}}} (\widehat{\mathbb{Z}}^{3}/\widehat{\mathbb{Z}}) ,
\]
où $\mu^{\leftarrow}_\infty = \varprojlim \mu_n$. Ainsi $C_\infty^{*}$ devient revêtement principal de $C_1^{*}$, de groupe $\Pi_\infty$.
\note{avant le dernier membre de (6), une première écriture $\simeq$ \struck{$\mathrm{T}$} $\mu^{\leftarrow}_\infty \otimes (\mathbb{Z}^{3}/\mathbb{Z})$ ; le numéro (4) ne figure pas sur la page, qui passe de (3'') à (5)}
\marginal{NB on laisse tomber l'indice $\mathbb{Q}$}

\page{26}
\textbf{1.3} Notons que $\mathfrak{S}_3$ opère sur $C_n$ de façon évidente, et les $C_{n'} \to C_n$ sont des $\mathfrak{S}_3$-morphismes, donc $\mathfrak{S}_3$ opère sur $C_\infty$. L'opération de $\mathfrak{S}_3$ sur $C_1$ ($= \mathbb{P}^{1}_{\mathbb{Q}}$) est celle déduite des opérations de $\mathfrak{S}_3$ sur $\{0, 1, \infty\}$. Enfin $\mathfrak{S}_3$ opère sur $\Pi_n$, de façon que $\mathfrak{S}_3$ opère sur $(C_n, \Pi_n)$ \ldots{} \add{donc $\mathfrak{S}_3$ opère sur $(C_\infty, \Pi_\infty)$\ldots}

\textbf{1.4} Considérons la clôture alg.\ $\overline{\mathbb{Q}}$ de $\mathbb{Q}$ dans $\mathbb{C}$, d'où par changement de base les schémas $\overline{C}_n$, $\overline{C}_n^{*}$, $\overline{C}_\infty$, $\overline{C}_\infty^{*}$, sur $\overline{C}_1 = \mathbb{P}^{1}_{\overline{\mathbb{Q}}}$ et $\overline{C}_1^{*} = \mathbb{P}^{1}_{\overline{\mathbb{Q}}} \setminus \{0, 1, \infty\}$. Sur les schémas $\overline{C}_n$ \struck{\ill{}}
\[
  \text{(7)} \qquad \Gamma = \mathrm{Gal}(\overline{\mathbb{Q}}/\mathbb{Q})
\]
opère de façon évidente, et aussi sur les groupes
\[
  \text{(8)} \qquad \overline{\Pi}_n = \Pi_n(\overline{\mathbb{Q}}) = \mu_n^{3}/\mu_n \simeq \mu_n \otimes_{\mathbb{Z}} \mathbb{Z}^{3}/\mathbb{Z} \simeq \mu_n \otimes_{\widehat{\mathbb{Z}}} (\widehat{\mathbb{Z}}^{3}/\widehat{\mathbb{Z}})
\]
\note{au-dessus de $\overline{\Pi}_n$, un ajout lu « $\Pi_n = ?$ » ; sous (8), une ligne biffée commençant par « où $\overline{\mu}_n = \mu_n(\overline{\mathbb{Q}})$ »}
et $\mathfrak{S}_3$. On voit tout de suite que ces opérations s'ajustent de façon à ce qu'on trouve le groupe
\[
  \text{(9)} \qquad \mathcal{E}_n = (\mathfrak{S}_3 \times \Gamma) \cdot_{1/2} \Pi_n
\]
\note{« $\cdot_{1/2}$ » rend sa notation « $\cdot_{\frac12}$ » placée en indice du point : produit semi-direct ; le numéro (9) est écrit sur un (8)}

\page{27}
où $\mathfrak{S}_3 \times \Gamma$ opère sur $\Pi_n$ via l'opération évidente de $\mathfrak{S}_3$ sur $\mathbb{Z}^{3}/\mathbb{Z}$, et de $\Gamma$ sur $\mu_n$ via le caractère fondamental
\[
  \text{(10)} \qquad \Gamma \xrightarrow{\ \chi_n\ } (\mathbb{Z}/n\mathbb{Z})^{*} .
\]
Passant à la limite, on trouve une opération de $\mathfrak{S}_3 \times \Gamma$ sur
\[
  \text{(11)} \qquad
  \begin{cases}
    \Pi_\infty = \mathrm{T}^{3}/\mathrm{T} \simeq \mathrm{T} \otimes \mathbb{Z}^{3}/\mathbb{Z} \simeq \mathrm{T} \otimes \widehat{\mathbb{Z}}^{3}/\widehat{\mathbb{Z}}, \\
    \text{où } \mathrm{T} = \varprojlim \mu_n = \mu^{\leftarrow}_\infty(\overline{\mathbb{Q}})
  \end{cases}
\]
où $\mathfrak{S}_3$ opère via son op.\ sur $\mathbb{Z}^{3} \to \widehat{\mathbb{Z}}^{3}$, et $\Gamma$ par le caractère fondamental
\[
  \text{(12)} \qquad \Gamma \xrightarrow{\ \chi_\infty = \chi\ } \widehat{\mathbb{Z}}^{*} .
\]
Bien entendu,
\[
  \text{(13)} \qquad
  \begin{cases}
    \mathfrak{S}_3 \times \Gamma = \mathrm{Aut}_{\mathbb{Q}}\, C_1 \\
    (\mathfrak{S}_3 \times \Gamma) \cdot_{1/2} \mu_n \overset{\mathrm{df}}{=} \mathcal{E}_n \simeq \mathrm{Aut}_{\mathbb{Q}}(C_n) \\
    (\mathfrak{S}_3 \times \Gamma) \cdot_{1/2} \mu_\infty \overset{\mathrm{df}}{=} \mathcal{E}_\infty \simeq \mathrm{Aut}_{\mathbb{Q}}(C_\infty)
  \end{cases}
\]

\textbf{1.5} En fait, l'\uncertain{hom.}\ \ill{}
\[
  \pi_1(\overline{C}_1^{*}) \longrightarrow \Pi_\infty
\]
\note{sous $\pi_1(\overline{C}_1^{*})$, écrit en colonne : $\Vert$ $\pi_1(\mathbb{P}^{1}_{\overline{\mathbb{Q}}} \setminus \{0,1,\infty\})$ ; sous $\Pi_\infty$ : $\simeq \mathrm{T}^{3}/\mathrm{T} \simeq \mathrm{T} \otimes \mathbb{Z}^{3}/\mathbb{Z}$}
définit un isom.

\page{28}
\[
  \text{(14)} \qquad \pi_1(\overline{C}_1^{*})^{\mathrm{ab}} \simeq \Pi_\infty \quad (\simeq H_1(\overline{C}_1^{*}))
\]
(i.e.\ on est en train de faire l'étude arithmétique des rev.\ \add{étales} abéliens de $\overline{C}_1^{*}$ \ldots).

\textbf{1.6} Considérons les fibres $(\overline{C}_\infty)_\zeta$, $\zeta = 0, 1, \infty$, disons \struck{$E_0$}, \struck{$E_1$}, $E_\infty$. Ainsi $E_0 = \varprojlim E_0(n)$,
\[
  \begin{aligned}
    E_0(n) &= \{(x, y, z) \in \overline{\mathbb{Q}}^{3} \mid x = 0,\ y^{n} + z^{n} = 0\}/\mu_n \\
    &\simeq \{(y, z) \in \overline{\mathbb{Q}}^{2} \mid (y/z)^{n} = -1\}/\mu_n \\
    &\simeq \{\lambda \in \overline{\mathbb{Q}} \mid \lambda^{n} = -1\}
  \end{aligned}
\]
$\simeq$ \struck{$\overline{\mathbb{Q}}$} image inverse de $-1 \in \mu_2$ par $\mu_{2n} \to \mu_2$, donc
\[
  \text{(15)} \qquad E_0 \overset{\varphi_0}{\simeq} \text{image inverse de } -1 \in \mu_2 \text{ par } \mu^{\leftarrow}_{\infty} \to \mu_2 .
\]
C'est pareil pour $\zeta = 0, 1, \infty$ (par $\mathfrak{S}_3$) et compatible \struck{de façon plus intrinsèque} avec l'action de \ill{} $\Gamma$, \struck{Pour formuler \ill{}}
\struck{de $\Pi_\infty$ opérant sur les deuxièmes \ill{} via l'action de $\mathrm{T}^{3}/\mathrm{pr}_\zeta + \mathrm{diag}(\mathrm{T})$}
de façon simple les compatibilités avec les actions de $\mathfrak{S}_3$ et $\Pi_\infty$, on va présenter
\note{la fin de la page est un bloc biffé de traits obliques, lisible par endroits seulement ; l'exposant de $\mu^{\leftarrow}$ dans (15) est incertain}

\page{29}
les choses de façon plus intrinsèque. Considérons l'homomorphisme canonique
\[
  \text{(16)} \qquad \Pi_\infty \longrightarrow \Pi_2
\]
et soit
\[
  \text{(17)} \qquad \Pi_2^{*} = \Pi_2 \setminus \{1\} \simeq (\mathbb{F}_2^{3}/\mathbb{F}_2) \setminus \{0\} .
\]
C'est isom.\ à l'ens.\ $I = \{0, 1, \infty\}$ de façon compatible avec l'action de $\mathfrak{S}_3$ [NB $\Gamma$ opère trivialement sur $\Pi_2^{*}$ comme sur $I = \{0, 1, \infty\}$]. Ceci dit \struck{pour \ill{}}, \struck{considérons}

\struck{(18') $I = \{0,1,\infty\}$, $E(I) ={}$ $\overline{C}_\infty(\overline{\mathbb{Q}})\,|\,I$ $= \coprod_{\zeta \in I} \overline{C}_{\infty,\zeta}$ $= E_0 \sqcup E_1 \sqcup E_\infty$ ; posant $\widehat{\mu}_\infty ={}$ $\mu_\infty \cup \{0\}$ :}
\struck{(19) $E(I) ={}$ $\{(\zeta_0, \zeta_1, \zeta_\infty) \in{}$ $\widehat{\mu}_\infty^{3} \mid$ un des $\zeta_i$ est nul$\}$}
\note{les deux lignes précédentes sont encadrées et biffées ensemble}

considérons la suite \struck{exacte} \uncertain{cx.}\ et \struck{$\mu_\infty$} de $(\mathfrak{S}_3 \times \Gamma)$-modules
\[
  \text{(18)} \qquad 0 \longrightarrow \Pi_\infty \xrightarrow{\ 2\ } \Pi_\infty \xrightarrow{\ p\ } \Pi_2 \longrightarrow 0, \qquad \Pi_2 \supset I = \{0, 1, \infty\},
\]
et pour $i \in I$ soit
\[
  \text{(19)} \qquad Z_i \subset \Pi_\infty \ (\simeq \mathrm{T})
\]
le sous-groupe de ramification correspondant, image dans $\Pi_\infty \simeq \mathrm{T}^{3}/\mathrm{T}$ du \struck{$i$-ème} facteur $\mathrm{T}$ d'indice $i$ dans $\mathrm{T}^{3}$. \struck{Soit enfin} On a donc
\[
  \text{(20)} \qquad E_i \simeq p^{-1}(\{i\})/Z_i .
\]
\marginal{Autre façon : on utilise les trois points rat.\ sur $\mathbb{Q}$, (19) des $\overline{C}_\infty$, \struck{d'ailleurs} \ill{} des $E_i$ ($i = 0, 1, \infty$), d'où $E_i \simeq \Pi_\infty/Z_i \simeq \mathrm{T}_{A}(\pm 1)(I \setminus \{i\})$ (20)}
\note{note écrite en oblique dans le coin inférieur gauche ; lecture d’ensemble incertaine, en particulier l’indice $A$ de $\mathrm{T}_{A}$}

\page{30}
\textbf{1.7} Soient
\[
  \text{(21)} \qquad \alpha = -j, \qquad \overline{\alpha} = -\overline{j}
\]
\marginal{NB $j = -\tfrac12 + i\tfrac{\sqrt3}{2}$, \ $-\overline{j} = \tfrac12 + i\tfrac{\sqrt3}{2}$ ; posons $\alpha_{*} = \{\alpha, \overline{\alpha}\}$, $E(\alpha_{*}) = E(\{\alpha, \overline{\alpha}\})$}
les deux pts de $\overline{C}_1$ invariants par $\mathfrak{S}_3^{+}$, i.e.\ deux éléments \add{en} fixes sous $\mathfrak{S}_3$ \uncertain{de card.\ 2}. Une \ill{} d'ensembles
\[
  \text{(22)} \qquad E(\alpha) \sqcup E(\overline{\alpha})
\]
image inverse de $\{\alpha, \overline{\alpha}\}$ en tant que $\mathcal{E}$-ensemble --- i.e.\ avec les op.\ de $\Gamma$, $\mathfrak{S}_3$, $\Pi_\infty$. On trouve, par calcul en coordonnées affines,
\[
  \text{(23)} \qquad
  \begin{cases}
    E(\alpha)_n = \{(\xi, \eta) \in \overline{\mathbb{Q}}^{2} \mid \xi^{n} = \overline{j},\ \eta^{n} = j\} \\
    E(\overline{\alpha})_n = \{(\xi, \eta) \mid \xi^{n} = j,\ \eta^{n} = \overline{j}\}
  \end{cases}
\]
\note{dans (23), un « $-1$ » biffé devant $j$}
\struck{Notons qu'alors les conditions \ill{} $\xi^{3n} = \eta^{3n} = 1$ ; donc $E(\alpha) = \varprojlim E(\alpha)_n$ $\leftarrow \varprojlim \mu_{3n} \times \mu_{3n}$ $\simeq \mu_\infty \times \mu_\infty$ ; $(\xi_n, \eta_n) \mapsto (\xi^{3}, \eta^{3})$}

Pour avoir une description intrinsèque des opérations de $\mathfrak{S}_3$ et de $\Pi_n$ (celles de $\Gamma$ sont données) on va \struck{écrire} plutôt \struck{\ill{}} prendre nos courbes affines projetantes de $C_n$, $C_1$, \struck{et prendre les points au-dessus de pt de $E_i(\alpha, \overline{\alpha}) \to$} $E_n(\{\alpha, \overline{\alpha}\})$, $E_1(\{\alpha, \overline{\alpha}\})$, \struck{\ill{}} $\widetilde{E}_n(\{\alpha, \overline{\alpha}\})$, $\widetilde{E}_1(\{\alpha, \overline{\alpha}\})$ --- puis \struck{\ill{}} passer au quotient pour avoir $E_n$, $E_1$ \ldots
\note{la seconde moitié de la page porte plusieurs lignes biffées en oblique, dont seules les formules ci-dessus se lisent}

\page{31}
\struck{$\widetilde{E}_1\{\alpha, \alpha\}$}

Soit \struck{$\widetilde{C}_1 = \{(\zeta_1, \zeta_2, \zeta_\infty) \in \mu_3^{3}$}
\[
  \text{(24)} \qquad \widetilde{C}_1 = \mathrm{Isom}(\underbrace{\{1, 2, \infty\}}_{I}, \mu_3) \simeq \{(\zeta_1, \zeta_2, \zeta_3) \in \mu_3^{3} \mid \zeta_1 \neq \zeta_2 \neq \zeta_3 \neq \zeta_1\}
\]
$\subset$ \struck{\ill{}} cône projetant de $\overline{C}_1 \subset \overline{\mathbb{Q}}^{3}$.
\note{le numéro de (24) est lu d'après la suite (25), (26) ; l'écriture est effacée}

$\widetilde{C}_1$ a $3! = 6$ éléments, et $\mu_3$ y opère diagonalement (multiplications scalaires), \struck{sans points fixes} librement --- le quotient $\widetilde{C}_1/\mu_3$ a donc deux éléments. L'application
\[
  \widetilde{C}_1/\mu_3 \longrightarrow \overline{C}_1
\]
\marginal{NB $\widetilde{C}_1/\mu_3 \subset \Pi_3$ est formé des \struck{\ill{}}}
devient une bijection
\[
  \text{(25)} \qquad \widetilde{C}_1/\mu_3 \xrightarrow{\ \sim\ } \overline{C}_1^{\,\mathfrak{S}_3^{+}} = \{\alpha, \overline{\alpha}\}
\]
mettant en évidence l'opération de $\mathcal{E}$ sur $(\alpha, \overline{\alpha})$, via ses opérations sur $\widetilde{C}_1 \subset \mu_3^{3} = \mu_3^{I}$. [NB $\mathcal{E}$ opère via $\mathfrak{S}_3 \times \{\pm 1\}$, via $\Gamma \xrightarrow{\chi_3} \{\pm 1\}$).

Soit alors $\widetilde{C}_n$ l'image inverse de $\widetilde{C}_1$ par l'application \uncertain{naturelle} des cônes projetants de $\overline{C}_n$, sur le cône projetant de $\overline{C}_1$, donc \struck{l'ensemble des $(x, y, z)$ \ill{} $\overline{C}_\infty$}
\[
  \text{(26)} \qquad
  \begin{aligned}
    \widetilde{C}_n &= \{(x, y, z) \in \overline{\mathbb{Q}}^{3} \mid (x^{n}, y^{n}, z^{n}) \in \widetilde{C}_1\} \\
    &= \{(x, y, z) \in \overline{\mathbb{Q}}^{3} \mid x^{3n} = y^{3n} = z^{3n} = 1,\ x^{n} \neq y^{n} \neq z^{n} \neq x^{n}\}
  \end{aligned}
\]
\note{dans la première ligne de (26), « $x^n + y^n + z^n = 1$, $x$ » biffé avant la condition}
Notons que $\mu_{3n}$ opère diagonalement sur $\overline{C}_n$,

\page{32}
L'application
\[
  \widetilde{C}_n \longrightarrow \widetilde{C}_1
\]
est compatible avec
\[
  \mu_{3n} \xrightarrow{\ \zeta \mapsto \zeta^{n}\ } \mu_n
\]
et induit donc par passage au quotient
\[
  \widetilde{C}_n/\mu_{3n} \longrightarrow \widetilde{C}_1/\mu_n .
\]
\note{ici et dans (27) il écrit $\mu_n$ là où $\widetilde{C}_1/\mu_3$ était défini, et où $\zeta\mapsto\zeta^n$ envoie $\mu_{3n}$ sur $\mu_3$ ; transcrit tel quel}
Ceci dit, on a un diagramme commutatif
\note{(27) ; flèche horizontale du bas : « projection canonique induite par $\overline{C}_n \to \overline{C}_1$ » ; sous chaque $\{\alpha,\overline{\alpha}\}$ une accolade marquée $\alpha_{*}$}

\begin{tikzcd}[column sep=large]
\widetilde{C}_n/\mu_{3n} \arrow[r] \arrow[d, "\wr"'] & \widetilde{C}_1/\mu_n \arrow[d, "\wr"] \\
E_n(\{\alpha, \overline{\alpha}\}) \arrow[r] & E_1(\{\alpha, \overline{\alpha}\})
\end{tikzcd}

Passant à la limite, on pose
\[
  \text{(28)} \qquad \widetilde{C}_\infty = \{(x, y, z) \in \mathrm{T}^{3} \ (= \mu_\infty^{3}) \mid x_3 \neq y_3 \neq z_3 \neq x_3\}
\]
où pour $\xi \in \mathrm{T}$, $\xi_3$ désigne son image par $\mathrm{T} \to \mathrm{T}/3\mathrm{T} \simeq \mu_3$. Donc
\[
  \widetilde{C}_\infty \subset \mathrm{T}^{3} \quad \text{stable par } \mathrm{T}
\]
et on trouve un \ill{}
\[
  \text{(29)} \qquad \widetilde{C}_\infty/\mathrm{T} \hookrightarrow \mathrm{T}^{3}/\mathrm{T} \simeq \Pi_\infty ,
\]
\note{en pied de page, coupé : $E^{21}(\alpha_{*})$ \ill{}}

\page{33}
de façon précise
\[
  \text{(30)} \qquad \widetilde{C}_\infty/\mu_\infty \simeq \text{image inverse de } \overbrace{\widetilde{C}_1/\mu_3}^{\omega} \subset \mu_3^{3}/\mu_3 = \Pi_3 \ \text{par } \Pi_\infty \to \Pi_3 .
\]
En termes de la suite exacte
\[
  \text{(31)} \qquad 0 \longrightarrow \Pi_\infty \xrightarrow{\ 3\ } \Pi_\infty \xrightarrow{\ p_3\ } \Pi_3 \longrightarrow 0, \qquad \omega \subset \Pi_3,
\]
$\Pi_\infty$ opère librement sur cette partie de $\Pi_\infty$, image inverse de $\omega$ par $\Pi_\infty \to \Pi_3$. Ceci donne les opérations de $\Pi_\infty$ sur $E(\alpha_{*}) \simeq p_3^{-1}(\omega)$, celles de $\Gamma$ et $\mathfrak{S}_3$ étant induites par leurs opérations sur $\Pi_\infty$\ldots

NB Sur $\Pi_3$, \struck{$\mathcal{E} \times$} $\mathfrak{S}_3 \times \Gamma$ opère via $\mathfrak{S}_3 \times \{\pm 1\}$ \add{(opérations \uncertain{non fidèles} sur $\Pi_3$)}, ($\Gamma \xrightarrow{\chi_3} \{\pm 1\}$). Considérons $\mathfrak{S}_3^{+}$ ($\simeq \mathbb{Z}/3\mathbb{Z}$) opérant sur $\Pi_3$ ($\simeq \mathbb{F}_3^{2}$), alors on sait qu'il doit y avoir une droite invariante --- \struck{et les \ill{}} en fait elle est unique --- et $\omega$ est formé des deux éléments non nuls de cette droite\ldots

\page{34}
\textbf{(1.8)} Considérons enfin les pts fixes
\[
  \text{(32)} \qquad \alpha'(0) = 2, \quad \alpha'(1) = -1, \quad \alpha'(\infty) = \tfrac12
\]
qui sont les pts fixes, autres que $0, 1, \infty$ resp., des transpositions $\tau_0, \tau_1, \tau_\infty \in \mathfrak{S}_3$.
\note{dans (32), les $\alpha'(i)$ sont écrits au-dessus de lettres biffées ($\beta_0$, $\beta_1$, $\beta_\infty$) ; même correction de $\beta$ en $\alpha'$ dans (33) et dans le dessin}
\drawing{un cercle (la droite projective réelle) portant, dans l'ordre, les points $\infty$, $\alpha'(0)$, $1$, $\alpha'(\infty)$, $0$, $\alpha'(1)$ ; chaque $\alpha'(i)$ est sur l'arc opposé au point $i$}
Soit
\[
  \text{(33)} \qquad \alpha'_{*} = \{\alpha'(0), \alpha'(1), \alpha'(2)\} .
\]
\note{« $\alpha'(2)$ » sic, pour $\alpha'(\infty)$}
On va déterminer $E(\alpha'_{*})$ \add{$= \coprod_{i \in I} E(\alpha'(i))$, comme un} $\mathcal{E}$-ensemble.

On trouve
\[
  \text{(34)} \qquad
  \begin{cases}
    E_n(\alpha'(i)) \simeq \{(\xi_j) \in \overline{\mathbb{Q}}^{I \setminus \{i\}} \mid \xi_j^{n} = -\tfrac12 \ \ \forall j \in I \setminus \{i\}\} \simeq K_n(-\tfrac12)^{I \setminus \{i\}} \\
    E(\alpha'(i)) \simeq K(-\tfrac12)^{I \setminus \{i\}}
  \end{cases}
\]
où $\forall g \in \mathbb{Q}$, on pose
\[
  \text{(35)} \qquad
  \begin{cases}
    K_n(g) = \{\xi \in \overline{\mathbb{Q}} \mid \xi^{n} = g\} \\
    K_\infty(g) = K(g) = \varprojlim_{n} K_n(g)
  \end{cases}
\]
$K_\infty(g)$ étant un torseur sous $\mathrm{T} = \mu_\infty$, sur lequel $\Gamma$ opère \uncertain{donc} (\uncertain{un} \struck{$\mathrm{Ext}$}$(\mu_\infty, \Gamma)$-torseur\ldots)

Ceci en \uncertain{ayant}, on \uncertain{voit} les \uncertain{iso.\ can.}
\[
  K(gg') \simeq K(g) \wedge_{\mathrm{T}} K(g')
\]
\[
  \text{(36)} \qquad K(-\tfrac12) \simeq K(-1) \wedge_{\mathrm{T}} K(\tfrac12) \simeq K(-1) \wedge_{\mathrm{T}} K(2)^{-1}
\]
\[
  \text{(37)} \qquad K(-1) \hookrightarrow \mathrm{T}, \qquad (1 \to \mathrm{T} \xrightarrow{\ 2\ } \mathrm{T} \xrightarrow{\ p_2\ } \mu_2 \to 1)
\]
\note{sous la flèche de (37) : « $\simeq p_2^{-1}(-1)$ » ; un symbole biffé devant $\mathrm{T}$}

\page{35}
Soit
\[
  \text{(38)} \qquad A = \mathrm{Aff}(\mathrm{T}, K(\tfrac12))
\]
\note{à la suite de (38), « $\simeq \mathrm{Aff}$ » biffé}
donc
\[
  \text{(39)} \qquad 1 \longrightarrow \mathrm{T} \longrightarrow A \longrightarrow \widehat{\mathbb{Z}}^{*} \longrightarrow 1 \qquad (\widehat{\mathbb{Z}}^{*} \simeq \mathrm{Aut}(\mathrm{T})) .
\]
\struck{Les ensembles ($\Pi_n$, $\Pi_\infty$, et les ensembles $E(I)$, $E(I)$ \add{groupes $\mu_n$, $\mu_\infty = \mathrm{T}$,} \ill{} \ill{}} \add{$E(\alpha_{*})$, $E(\alpha_{*})$}
$A$ opère sur $\Pi_n$, $\Pi_\infty$ \add{$= \mathrm{T}^{3}/\mathrm{T}$} par l'intermédiaire de $\widehat{\mathbb{Z}}^{*}$, \struck{$\Pi \to A$} (par homothéties \uncertain{\ldots}) et le groupe
\[
  \widetilde{\mathcal{E}} = (\mathfrak{S}_3 \times A) \cdot_{1/2} \Pi_\infty
\]
opère \add{de façon évidente} sur les ensembles $E(\alpha_{*})$ \struck{\ill{}} $E(\alpha_{*})$, $E(\alpha_{*})$, $E(\alpha'_{*})$, $E(\alpha'_{*})$ --- en fait les deux premiers opèrent via le quotient \struck{$\widetilde{\mathcal{E}}/\mathrm{T} \simeq$}
\[
  \widetilde{\mathcal{E}}/\mathrm{T} \simeq (\mathfrak{S}_3 \times \widehat{\mathbb{Z}}^{*}) \cdot_{1/2} \Pi_\infty .
\]
\marginal{$\alpha(0) = 0$, $\alpha(1) = 1$, $\alpha(\infty) = \infty$, $\alpha(*) = \{\alpha(0), \alpha(1), \alpha(\infty)\}$}
\note{la liste des ensembles répète $E(\alpha_{*})$ et $E(\alpha'_{*})$ : sans doute une fois soulignés (ensembles sur $\overline{\mathbb{Q}}$), une fois non ; la note de marge ajoute $\alpha_{*} = \{0,1,\infty\}$ à côté du $\{\alpha,\overline{\alpha}\}$ de 1.7}

D'autre part, on a
\[
  \text{(40)} \qquad \Gamma \longrightarrow A
\]
\uncertain{homom.}\ dont l'image $\Gamma'$ \uncertain{est} d'indice fini (\uncertain{est-ce} $=$ surjectif ?) \marginal{?} $\Gamma'$ est le groupe de Galois de la sous-extension de $\overline{\mathbb{Q}}/\mathbb{Q}$ engendrée par les racines \uncertain{$n$-ièmes} de $\tfrac12$,

\page{36}
pour $n$ variable. Ceci dit, on a un homom.\ évident
\[
  \text{(41)} \qquad \mathcal{E} = (\mathfrak{S}_3 \times \Gamma) \cdot_{1/2} \Pi_\infty \longrightarrow \widetilde{\mathcal{E}} = (\mathfrak{S}_3 \times A) \cdot_{1/2} \Pi_\infty
\]
déduit de (40), et les opérations de $\mathcal{E}$, sur les ens.\ précédents se font \uncertain{à travers} $\widetilde{\mathcal{E}}$\ldots

\textbf{(1.9)} Les développements précédents n'utilisaient pas que \struck{$\mathbb{Q}$} la clôture alg.\ $\overline{\mathbb{Q}}$ \add{\uncertain{envisagée}} de $\overline{\mathbb{Q}}$ était plongée dans $\mathbb{C}$ --- ils \uncertain{resteraient valables} en remplaçant $\mathbb{Q}$ par n'importe quel corps premier, à condition de \ill{} \ill{} \ill{} premiers à la car. ($\overline{C}_\infty$ est alors le revêtement universel abélien \emph{modéré} de $(\overline{C}_1, I)$\ldots). Il faut l'utiliser, pour faire le lien avec les cartes \uncertain{3-gradiées} finies abéliennes, et le \uncertain{point} \uncertain{cartographique}. On va d'abord trivialiser $\overline{C}_\infty$ au-dessus de $\alpha = -j \in \omega = \{\alpha, \overline{\alpha}\}$, pour pouvoir l'identifier au rev.\ universel profini \struck{pointé} \emph{en $\alpha$}.

\page{37}
Rappelons
\[
  \text{(42)} \qquad E(\alpha) \subset \Pi_\infty \xrightarrow{\ p_3\ } \Pi_3, \qquad E(\alpha) = p_3^{-1}(\alpha)
\]
\[
  \alpha = (\overline{j}, j, 1) \bmod \mu_3 = \{(\overline{j}, j, 1), (j, 1, \overline{j}), (1, \overline{j}, j)\} \in \Pi_3
\]
\note{les deuxièmes et troisièmes coordonnées portent des surcharges ; sous la ligne, un mot biffé}
\struck{Posons.} Nous \struck{prendrons} utilisons
\[
  \Bigl( \quad j = \exp\frac{2i\pi}{3}, \qquad \overline{j} = \exp -\frac{2i\pi}{3}
\]
et posons
\[
  \text{(43)} \qquad \zeta_n = \exp\frac{2i\pi}{3n}
\]
et
\[
  \begin{cases}
    \alpha_n = (\overline{\zeta}_n, \zeta_n, 1) \\
    \alpha_\infty = (\overline{\zeta}_\infty, \zeta_\infty, 1)
  \end{cases}
\]
\note{les deux lignes sont précédées et surchargées de ratures ; dans $\alpha_n$, la deuxième coordonnée porte un « $j$ » biffé ; $\alpha_\infty$ est réécrit sous une première forme entièrement raturée}
\marginal{[pour $\alpha_n$] considérés sur \uncertain{comme} un élément de $\overline{C}_n$ au-dessus de $\alpha = -j$ \struck{mod $\mu_n$ dans} sont comme élément de $\mu_{3n}^{3}/\mu_{3n}$ \struck{\ill{}} $= \Pi_{3n}$ ; [pour $\alpha_\infty$] considérés sur \uncertain{comme} pts de $\overline{C}_\infty$ au-dessus de $\alpha = -j$ sont comme él.\ de $\mu_\infty^{3}/\mu_\infty = \Pi_\infty$}
où
\[
  \text{(44)} \qquad \zeta_\infty = (\zeta_n) \in \mathrm{T}
\]
est le générateur canonique de $\mathrm{T}$, donné par l'exponentielle complexe, permettant donc de définir un iso.\ can.
\[
  \text{(45)} \qquad \widehat{\mathbb{Z}} \overset{\varphi}{\simeq} \mathrm{T} .
\]

\page{38}
Maintenant on déduit de là des isom.
\[
  \text{(46)} \qquad
  \begin{cases}
    \Pi_\infty \simeq \widehat{\mathbb{Z}}^{3}/\widehat{\mathbb{Z}} \\
    \mathcal{E} \simeq (\mathfrak{S}_3 \times \Gamma) \cdot_{1/2} (\widehat{\mathbb{Z}}^{3}/\widehat{\mathbb{Z}})
  \end{cases}
\]
\note{entre les deux lignes de (46), une ligne biffée, lue « $\mathcal{E}$ $(\mathfrak{S}_3 \times \widehat{\mathbb{Z}}^{*})/\widehat{\mathbb{Z}}$ » ; sous $\mathfrak{S}_3 \times \Gamma$ : « opère via permutation des facteurs et via homothétie par $\chi : \Gamma \to \widehat{\mathbb{Z}}^{*}$ »}
\[
  \Pi_n \simeq (\mathbb{Z}/n)^{3}/(\mathbb{Z}/n)
\]
\note{cette formule est entourée, et une flèche la renvoie à la première ligne de (46)}
\[
  \text{(47)} \qquad \alpha_\infty \simeq (-1, +1, 0) \bmod \widehat{\mathbb{Z}}
\]
d'où
\[
  \text{(47')} \qquad \alpha'_\infty = \overline{\alpha}_\infty = (1, -1, 0) \bmod \widehat{\mathbb{Z}} .
\]
Considérons les éléments
\[
  \text{(48)} \qquad
  \begin{cases}
    \rho_0, \rho_1, \rho_\infty \in \Pi_\infty \\
    \rho_0 = (1, 0, 0) \bmod \widehat{\mathbb{Z}} \\
    \rho_1 = (0, 1, 0) \bmod \widehat{\mathbb{Z}} \\
    \rho_\infty = (0, 0, 1) \bmod \widehat{\mathbb{Z}}
  \end{cases}
\]
on sait que
\[
  \text{(49)} \qquad \rho_0 \rho_1 \rho_\infty = 1 .
\]
\marginal{Ce sont ces éléments qui donnent l'isom.\ de $\widehat{\Gamma}^{+}_{12}$ avec $\Pi_\infty$.}
\note{la dernière phrase est marquée d'un double trait vertical dans la marge gauche ; l'indice de $\widehat{\Gamma}^{+}$ est lu « 12 » sans assurance}

\page{39}
\struck{Il reste à définir $b$}

Il faut expliciter $K(2)$, \struck{\ill{}} en choisissant une origine
\[
  \text{(50)} \qquad \theta_\infty = (\underbrace{\sqrt[n]{2}}_{\theta_n})_{n \in \mathbb{N}^{*}}
\]
donc $K(2)$ est identifié à $\mathrm{T}$ donc à $\widehat{\mathbb{Z}}$, et
\[
  \text{(51)} \qquad A \simeq \mathrm{Aff}(1, \widehat{\mathbb{Z}}) \simeq \widehat{\mathbb{Z}}^{*} \cdot_{1/2} \widehat{\mathbb{Z}}
\]
opérant sur $K(2) \simeq \widehat{\mathbb{Z}}$ de façon tautologique
\[
  \text{(52)} \qquad (a, b)(x) = ax + b \qquad (a \in \widehat{\mathbb{Z}}^{*},\ b \in \widehat{\mathbb{Z}},\ x \in K(2)) .
\]
On a donc
\[
  \text{(53)} \qquad K(\tfrac12) \simeq \widehat{\mathbb{Z}}
\]
mais opérations
\[
  \text{(54)} \qquad (a, b)(x) = ax - b \qquad (a \in \widehat{\mathbb{Z}}^{*},\ b \in \widehat{\mathbb{Z}},\ x \in K(\tfrac12)) .
\]
Ceci permet d'expliciter entièrement $E(\beta_{*})$ en plus de $E(\omega)$\ldots
\note{« $\beta_{*}$ » : il n'a pas partout remplacé l'ancienne lettre $\beta$ par $\alpha'$ (cf.\ p.~34)}

Il reste \add{enfin} : expliciter les \uncertain{bijections} de \uncertain{transport parallèle}
\[
  E(\alpha) \longrightarrow E_0,\ E_1,\ E_\infty
  \qquad \text{et} \qquad
  E(\alpha) \xrightarrow{\ \sim\ } E'_0 = E_2,\ \ E'_1 = E_{-1},\ \ E'_\infty \simeq E_{1/2}
\]
\note{deux faisceaux de trois flèches partant de $E(\alpha)$ ; ceux de droite marqués $\sim$}
[et $E'_0 \rightrightarrows E_1, E_\infty$ ; \ $E'_1 \rightrightarrows E_\infty, E_0$ \struck{\ill{}} ; \ $E'_\infty \rightrightarrows E_0, E_1$]
d'où on déduit les bijections correspondantes \marginal{?}
de $E(\overline{\alpha})$ en conjuguant par la conjugaison complexe. Comme on \uncertain{sait} \uncertain{a priori} des \uncertain{bijections}

\page{40}
compatibles avec l'action de $\Pi_\infty$ et de $\mathfrak{S}_3$ \struck{\ill{}} $\times \{1, \tau\}$ ($\tau$ étant la conjugaison complexe) il suffit pour déterminer ces 18 flèches, de déterminer deux d'entre elles, \struck{telles que} \add{p.\,ex.}
\[
  E(\alpha) \longrightarrow E'_\infty \longrightarrow E_0
\]
\note{sous $E(\alpha)$ : « $\Vert\ \alpha_1$ » ; sous $E'_\infty$ : « $\Vert\ E_{1/2}$ »}
[NB $E(\alpha) \to E_0$ est le composé] et $=$ de connaître \add{les} images dans $E_0$, $E_1$ de l'élément marqué $\alpha_\infty$ de $E(\alpha)$, [d'où élément marqué $\alpha'_\infty(i) \in E'_i = E(\alpha'(i))$, et $\alpha_\infty(i) \in E_i = E(\alpha(i))$, $\forall i \in \{0, 1, \infty\}$].

On trouve
\[
  \begin{aligned}
    \alpha'_n(\infty) &= \Bigl(\sqrt[n]{\tfrac12}\, \exp\frac{-\pi i}{n},\ \sqrt[n]{\tfrac12}\, \exp\frac{\pi i}{n},\ 1\Bigr) \\
    &= (\theta_n^{-1} \overline{\zeta}_n,\ \theta_n^{-1} \zeta_n,\ 1) \\
    \alpha'_\infty(\infty) &= (\theta_\infty \wedge \overline{\zeta}_\infty,\ \theta_\infty \wedge \zeta_\infty,\ 1)
  \end{aligned}
\]
\note{devant le second « $=$ », une écriture raturée ; dans la dernière ligne, le signe entre $\theta_\infty$ et $\zeta$ est lu $\wedge$ sans assurance, et l'exposant $-1$ de la ligne précédente n'y figure pas. Avec $\zeta_n = \exp\frac{2i\pi}{3n}$ de (43), $\exp\frac{\pi i}{n}$ n'est pas $\zeta_n$ : transcrit tel quel}
\note{la page s'arrête ici, au milieu de la détermination des flèches ; l'argument continue au-delà du lot}

\end{document}
